\documentclass[11pt]{article}
\usepackage{fontspec}
\setmainfont{Times New Roman}
\setsansfont{Arial}
\setmonofont{Consolas}
\usepackage{amsmath,amssymb,mathtools}
\usepackage{geometry}
\usepackage{microtype}
\geometry{a4paper,margin=2.35cm}
\setlength{\parindent}{1.25em}
\setlength{\parskip}{0pt}
\makeatletter
\renewcommand\footnotesize{\@setfontsize\footnotesize{7.7}{8.7}\selectfont}
\makeatother
\emergencystretch=100em
\allowdisplaybreaks
\sloppy
\hbadness=10000
\vbadness=10000
\hfuzz=2pt
\vfuzz=10pt
\raggedbottom

\begin{document}

% Unit: Paper31 Section06 no. 2 v001. Direct Interslavic Latin translation from the German final-audited source slice, source lines 346--350 / segment P31-S0076, checked against printed scan/OCR witness page 97. Footnote counter continues after Paper31 Section06 no. 1.
\setcounter{footnote}{27}

\textbf{2. \emph{Diskriminantne idealy primarnyh kolc s algebraično zamknjenym podkolcom $P$.}} Ako nulovy ideal kolca $\mathfrak R$ jest prosty ideal i $P$ jest algebraično zamknjeno, togda $\mathfrak R$ imaje rang jedan, zato jest identično s $P$ (\S~3, 3); jediničny element $e$ možno vzęti kako modulnu bazu. Tym dostava se
\[
        S_{(\mathfrak R)}(e^2)=S_{(\mathfrak R)}(e)=e,
\]
i diskriminantny ideal jest ravny jediničnomu idealu.

\end{document}
